\chapter{The Fischer-Paterson Algorithm}

\section{Generalized Linear Products}
As discussed earlier, classical linear-time algorithms fail because the wildcard symbol destroys the transitivity of character equivalence. Fischer and Paterson \cite{fp1974} recognized that to solve this problem, they had to abandon standard character-by-character comparisons entirely. 

To bypass this limitation, they redefined string matching using the concept of generalized linear products. Instead of sliding the pattern and comparing strings directly, the relationship between a text sequence $T$ and a pattern sequence $P$ is expressed as a convolution. The resulting match vector $Z$ is defined using custom mathematical operators $\otimes$ (representing character evaluation) and $\oplus$ (representing the summation of mismatches):
\begin{equation}
    Z_k = \bigoplus_{i+j=k} T_i \otimes P_j \quad \text{for} \quad k=0, \dots, m+n
\end{equation}

\section{Boolean Product Reduction}
The key idea behind the Fischer-Paterson approach is flipping the problem upside down: instead of actively verifying if characters match, the algorithm searches strictly for collisions (mismatches). 

The algorithm iterates over all pairs of distinct symbols from the alphabet (where $\sigma \neq \tau$). For each symbol, it creates a binary indicator sequence where positions containing that specific letter are marked as $1$, and all others (including wildcards) as $0$. 

If either comparing character is a wildcard, the indicator mask yields 0, negating the collision. Thus, exact collisions occur if and only if two differing standard characters are aligned. A valid alignment at a given shift means that zero collisions occurred between any two different letters. Using De Morgan's laws, this absence of collisions can be mathematically expressed as the negation of a Boolean product:
\begin{equation}
    \bigwedge_{i+j=k} T_i \equiv P_j \leftrightarrow \neg \bigvee_{i+j=k} \hat{T}_i \land \hat{P}_j
\end{equation}

\section{Polynomial Arithmetization}
Evaluating this Boolean product directly in memory is inefficient. To solve this, Fischer and Paterson translated the Boolean logic into standard arithmetic. This operation is mathematically identical to polynomial multiplication, which they implemented by packing the binary sequences into massive integers.

To ensure that adding these binary coefficients does not cause bits to overlap and corrupt adjacent values (overflow), they introduced a buffer constant $r$, chosen so that $2^r > n+1$. The binary vectors are then evaluated as polynomials at the point $2^r$, safely spacing the values apart:
\begin{equation}
    T(2^r) = \sum_{i=0}^{n-1} T_i \cdot 2^{ri}, \quad P(2^r) = \sum_{j=0}^{m-1} P_j \cdot 2^{rj}
\end{equation}

Algorithm 2 shows this bitwise arithmetization in practice. The buffer constant $r$ is calculated (line 3) to prevent integer overflow. The main loop (line 4) iterates over all distinct character pairs from the defined alphabet. The character sequences are packed into large integer representations using bit shifts. We invoke \textsc{Reverse} on the pattern to structurally transform standard multiplication into equivalent signal convolution. The numbers are multiplied using arbitrary-precision arithmetic. The relevant block overlaps generated by the multiplication are extracted via \textsc{ExtractBlocks} (which isolates the bits at position $k \cdot r$) and accumulated into the result array $Z$. If a position in $Z$ is exactly zero, it means no collisions occurred, and the match is perfectly valid.

\begin{algorithm}[H]
\caption{Fischer-Paterson Bitwise Multiplication}
\begin{algorithmic}[1]
\Function{FPBitwise}{$T$: string, $P$: string, $\Sigma$: set}
    \State $Z[0 \dots n-m] \gets 0$
    \State $r \gets \lceil \log_2(|T| + 1) \rceil$ \Comment{Bit-shift buffering constant}
    \For{\textbf{each} distinct pair $(\sigma, \tau) \in \Sigma \times \Sigma$}
        \State $X_{int} \gets \text{\textsc{PackToInteger}}(T, \sigma, r)$
        \State $Y_{int} \gets \text{\textsc{PackToInteger}}(\text{\textsc{Reverse}}(P), \tau, r)$
        \State $P_{mult} \gets X_{int} \cdot Y_{int}$ \Comment{Arbitrary-precision multiplication}
        \State $Z \gets Z + \text{\textsc{ExtractBlocks}}(P_{mult}, r)$
    \EndFor
    \State \textbf{return} $\{i \colon Z[i] = 0\}$
\EndFunction
\end{algorithmic}
\end{algorithm}

\section{Complexity Analysis}

\begin{theorem}[Turing Machine Bound \cite{fp1974}]
For a finite alphabet, the `don't care' product of strings of lengths $m$ and $n$, $(m \ge n)$, can be computed with a multitape Turing machine in time $\mathcal{O}(m \cdot (\log n)^2 \cdot \log \log n)$.
\end{theorem}

While this avoids the quadratic $\mathcal{O}(n \cdot m)$ trap of the naive approach, it introduces a major new downside in the RAM model. 

\begin{theorem}[RAM Model Complexity]
The practical time complexity of the Fischer-Paterson integer multiplication approach utilizing the Karatsuba algorithm \cite{karatsuba1962} in a standard RAM model is $\mathcal{O}(|\Sigma|^2 \cdot (m \log_2 n)^{1.585})$.
\end{theorem}

\begin{proof}
The Fischer-Paterson algorithm maps text and patterns into large integers to exploit the formal similarity between string matching with don't-cares and integer multiplication. By substituting their asymptotic multiplier with the Karatsuba algorithm \cite{karatsuba1962} for processing these large numbers, the complexity of multiplying two $N$-bit integers becomes $\mathcal{O}(N^{\log_2 3}) \approx \mathcal{O}(N^{1.585})$. Given that the integer sizes scale to $N \approx m \log_2 n$ bits, and the algorithm iterates over all distinct character pairs ($|\Sigma|^2$), the resulting total time complexity is $\mathcal{O}(|\Sigma|^2 \cdot (m \log_2 n)^{1.585})$. 
\end{proof}

Note that while Theorem 3.2 relies on Karatsuba's method, Harvey and van der Hoeven \cite{harvey2021} recently proved that $N$-bit integers can be multiplied in $\mathcal{O}(N \log N)$ time. Applying this bound directly lowers the asymptotic time complexity in the RAM model. 

However, regardless of this multiplication speedup, the massive packed integers still require dynamic memory allocation on the order of $\mathcal{O}(m \log n)$ bits. Combined with the $\mathcal{O}(|\Sigma|^2)$ dependency, this makes the algorithm impractical for large modern character sets like Unicode.

\section{Modern FFT Realization (Baseline Convolution)}
While the Fischer-Paterson paper relied on massive integers to compute the convolution, modern frameworks bypass this by operating directly in the frequency domain using the Fast Fourier Transform (FFT). 

This eliminates the need for manual bit-buffer management and overflow protection. More importantly, instead of searching for collisions between pairs of different characters ($\mathcal{O}(|\Sigma|^2)$), the FFT method simply calculates the positive matches for each character independently. This effectively reduces the alphabet dependency from $|\Sigma|^2$ down to $|\Sigma|$.

Algorithm 3 outlines this approach. First, the input vectors are padded to the nearest power of two $L$ (line 2) to optimize FFT performance. For each character $\sigma$ in the alphabet (line 4), two distinct boolean masks are generated (lines 5-6): an indicator mask for the text (locating occurrences of $\sigma$) and a mismatch mask for the reversed pattern (flagging characters different from $\sigma$, while safely ignoring wildcards). These masks are transformed into the frequency domain via forward FFTs (lines 7-8), multiplied point-wise, and subsequently transformed back using the Inverse FFT (line 9). By accumulating these results, the convolution counts the number of character collisions at each shift. Finally, the algorithm returns all indices $i$ where $Z[i] = 0$ (line 11), correctly identifying the shifts with zero mismatches.

\begin{algorithm}[H]
\caption{Baseline Character-Masked FFT}
\begin{algorithmic}[1]
\Function{BaselineFFT}{$T$: string, $P$: string, $\Sigma$: set}
    \State $L \gets 2^{\lceil \log_2(|T| + |P| - 1) \rceil}$ \Comment{Power-of-two padding length}
    \State $Z[0 \dots L-1] \gets 0$
    \For{\textbf{each} $\sigma \in \Sigma$}
        \State $U \gets \text{\textsc{IndicatorMask}}(T, \sigma, L)$
        \State $V \gets \text{\textsc{MismatchMask}}(\text{\textsc{Reverse}}(P), \sigma, L)$ 
        \State $\hat{U} \gets \text{\textsc{FFT}}(U)$
        \State $\hat{V} \gets \text{\textsc{FFT}}(V)$
        \State $Z \gets Z + \text{\textsc{IFFT}}(\hat{U} \odot \hat{V})$
    \EndFor
    \State \textbf{return} $\{i \colon Z[i] = 0\}$ \Comment{Return shifts with exactly 0 mismatches}
\EndFunction
\end{algorithmic}
\end{algorithm}

\begin{theorem}[Baseline FFT Complexity \cite{fp1974, cooley1965algorithm}]
The modern FFT realization of the Fischer-Paterson approach improves the overall time complexity to $\mathcal{O}(|\Sigma| \cdot n \log n)$.
\end{theorem}
\begin{proof}
Because the algorithm evaluates each unique character from the alphabet $\Sigma$ independently rather than evaluating cross-pairs $\Sigma \times \Sigma$, the loop executes exactly $|\Sigma|$ times. Inside the loop, the limiting operations are the forward and inverse Fast Fourier Transforms. Given a padded length $L \approx 2n$, each transform executes in $\mathcal{O}(n \log n)$ time, culminating in a total execution time of $\mathcal{O}(|\Sigma| \cdot n \log n)$.
\end{proof}

\section{Infinite Alphabets and Lower Bounds}
Because the execution time of the polynomial approach grows alongside the alphabet size, it forces an important theoretical question: is there an optimal algorithm that works independently of the alphabet by only using simple equality checks ($a = b$)?

Fischer and Paterson proved that if an algorithm is restricted strictly to basic equality comparisons, no method can beat the naive $\mathcal{O}(m \cdot n)$ bound. An intuitive adversarial argument proves this: if an algorithm skips even a single comparison involving a wildcard, an adversary could later replace that wildcard with a conflicting character, forcing the algorithm to fail. Thus, all pairs must be evaluated \cite{fp1974}.

However, if the machine is allowed to explicitly check if a character is specifically a wildcard (e.g., checking if $X_i$ is a wildcard), the rules of the problem change. The authors showed that by actively managing equivalence classes in memory—which allows tracking specific character sets to avoid rechecking existing wildcard constraints—the number of required comparisons can drop to linear $\mathcal{O}(m+n)$. Unfortunately, the computational overhead required to manage these memory structures is so high that, in practice, it remains slower than the simple naive algorithm.